To obtain the covector in Proposition~\ref{prop:residue-plane-neighbor-test},
we compare the actual pairing with its model on prescribed cut
subspaces. The corrected pairing estimate below uses the same original multidegree as the source lift in
Lemma~\ref{lem:lift-plane-symbol}.
\paragraph{The Fock metric and the signed cuts.}
\label{par:plane-gram}
\hypertarget{source-plane-gram}{}
We give the pairing argument in its signed polynomial model.
Normalize the valuation by \(v(p)=1\). For a monomial with
exponent \(\alpha\), write \(\alpha!=\prod_x\alpha_x!\) and
\(h_\alpha=v(\alpha!)/2\). Use the bilinear form and its
associated coefficient bound
\[
  [F,G]=\sum_\alpha F_\alpha G_\alpha\,\alpha!,
  \qquad \nu(F)=\min_\alpha \bigl(v(F_\alpha)+h_\alpha\bigr).
\]
Then \(v([F,G])\ge \nu(F)+\nu(G)\). In true multidegree
\(d=(d_i)\), set
\[
 n_0=\frac12 v\!\left(\prod_i d_i!\right),\qquad
 K_d=\sum_i r\bigl(1+\log_p(d_i+r)\bigr),\qquad
 e_d=r^2\bigl(1+\log_p(2\textstyle\sum_i d_i+r)\bigr).
\]

In this multidegree \(2n_0-K_d\le v(\alpha!)\le2n_0\).
Indeed the difference is the valuation of a product of
multinomial coefficients. Each multinomial is a product of
at most \(r\) binomial coefficients with upper entries at most
the letter degree. A binomial with upper entry \(n\ge1\) has
valuation at most \(1+\log_p n\), by the base-\(p\) carry
formula, which follows from Legendre's formula for factorials.
The same argument for a split \(\alpha=A+M\) gives
\[
   v(A!)+v(M!)\ge 2n_0-2K_d,
   \qquad \text{when \(A+M\) has multidegree \(d\)}.
\]

\begin{lemma}[Projectors for signed cuts in the factorial norm]\label{lem:plane-signed-projectors}
For every prescribed cut and threshold in fixed multidegree
\(d\), there is a Fock-self-adjoint rational projector onto
the signed cutoff subspace, whose loss in \(\nu\) is at most
\(e_d\). For nested or disjoint cuts write \(P_i\) for their projectors and \(P\) for the stipulated product. Then
\[
 P^2=P,\qquad \operatorname{im}P=\bigcap_i\operatorname{im}P_i,
\]
with loss at most \(e_d\) per cut.
It preserves multidegree and the monomial exponents outside
the union of those cuts.
\end{lemma}

\begin{proof}
\proofstep{The projector for a signed cut.}
Using the signed, generically conjugated Laurent order, the space attached to a cut \(S\) is
\[
 U(S,b)=\{F:\operatorname{ord}_S(F)\ge b\}.
\]
In fixed
multidegree it is the sum of the \(\mathrm{GL}_r\)-types
for the action on \(S\) whose least weight \(\lambda_r\) is
\(\ge b\). Indeed all torus exponents in a type are at least
\(\lambda_r\) in the selected coordinate, and the conjugates
of that torus test detect the lowest weight in every nonzero
copy. More explicitly the intersection of the conjugate
weight-cut spaces is invariant, and the copies it contains
are exactly the stated types. The type sum has a Fock
self-adjoint projector \(\Pi\), defined over \(\mathbb Q\).

By the trace formula and divided-power estimate for type
projections proved above, its matrix on the coefficient lattice satisfies
\[
 v_p(\Pi_{\beta\alpha})\ge-e_d\qquad(\text{all }\alpha,\beta).
\]
In applying the
estimate the coordinate span on the cut is at most \(\sum_i d_i\),
and it suffices to use \(2\sum_i d_i\) in the Hom weight bound.
Zero degree and \(r=1\) give scalar projections on the affected
monomial spaces. The trace formula works for the sum of types
with the same bound by the ultrametric inequality.

For the matrix entry \(\Pi_{\beta\alpha}\), self-adjointness
gives
\[
 \beta!\Pi_{\beta\alpha}=\alpha!\Pi_{\alpha\beta}.
\]

Consequently the two numbers
\(v(\Pi_{\beta\alpha})+h_\beta-h_\alpha\) and
\(v(\Pi_{\alpha\beta})+h_\alpha-h_\beta\) are equal.
Their common value is the mean of the two unweighted
entry valuations, each bounded below by \(-e_d\).
For the coefficient of \(x^\beta\) in \(\Pi F\) the
ultrametric inequality therefore gives
\[
 v((\Pi F)_\beta)+h_\beta
    \ge\min_\alpha\{v(F_\alpha)+h_\alpha-e_d\}.
\]
Taking the minimum in \(\beta\) yields the claimed norm bound.

For nested cuts the smaller-cut projector commutes with
the action on that cut and with the action on its complement.
It therefore preserves every isotypic component for the
larger cut and commutes with its type projector. For disjoint
cuts the projectors act on separate factors. A product over
such cuts is self-adjoint, fixes their intersection and loses
at most \(e_d\) per cut. It preserves multidegree and each
exponent outside the union of the cuts. This proves the projection and gives the product bound used below.

\end{proof}

\begin{paramdefinition}[The true and model lattices, cuts and signs]\label{hyp:plane-gram-data}
The block letters \(B_j\) in each frame \(j\) have specified
signs and total block degree \(b_j\). The shared letters \(X_t\),
indexed by \(t=(j,m)\), \(0\le m<r\), are vectors. Let
\(P_t=p^{g_t}\) be their radices. The model has the same
block letters and three vector letters \((L_t,W_t,T_t)\)
over each shared letter. Its collapse \(\phi\) sends them
respectively to \(X_t,X_t^{P_t},X_t\), coordinate by coordinate.
The exponents on \(L_t,W_t,T_t\) will be denoted by
\(\rho_t,\beta_t,M_t\). An admissible monomial has
\[
 \begin{gathered}
 0\le\rho_{t,a}<P_t, \\ M_{j,0}=0, \\ |M_{j,m}|\le b_j,
\end{gathered}
\]
and exact collapsed multidegree \(d\).

\smallskip
The cuts are defined as follows.
A true link cut contains \(B_j\) and the shared letters
\(X_{j,m'}\) for \(m'>m\). Denote it by \(A_{jm}\);
we use \(m+2\le r\). The corresponding model link cut
\(U_{jm}\) contains \(B_j\), the letters \(T_{j,m'}\)
for \(m'>m\), and all the \(W_{j-1,m'}\), where the frame
indices are cyclic.

A model chain cut \(S_{ji}\) contains a
nested subset of \(B_j\) and all the \(W_{j-1,m'}\).
Thus the link cuts alone, the chain cuts alone, and the
union of the chain cuts with the model link cuts are each
families of nested or disjoint cuts. Set
\[
 D_l=J(r-1)e_d,\qquad D_m=e_d\sum_j n_j,
\]
where \(J\ge1\) is the number of frames and \(n_j\) the
number of chain cuts at \(j\). The model total degree is
at most \(\sum_i d_i\) on admissible monomials because
each \(P_t\ge1\).

Use a diagonal unit \(\gamma(\alpha)\) on the true
monomials. Its two signs admit integral binomial polynomial
approximations of total degree at most \(6c\) modulo \(p^c\).

Suppose that \(\nu(F)=0\) in true multidegree \(d\), that
the true link orders of \(F\) are at least \(\ell_{jm}\),
and that each digit part of \(F\) has chain orders at least
\(\theta_{ji}\). The digit part for a fixed residue keeps the
block exponents and reads on \(W_t\) the quotient exponent
\[
 m_t=\rho_t+P_t\beta_t\quad\longmapsto\quad \beta_t=(m_t-\rho_t)/P_t,
\]
with no \(L_t,T_t\) exponents.
The pairing under consideration is \([F,\gamma\,\phi(Z)]\),
where \(\gamma\) acts coefficientwise. Use the inequalities
\begin{equation}\label{eq:plane-gram-reserves}
 \begin{split}
    c&>4(D_m+D_l+3K_d),\\
    s_{ji}&\ge r\, b_{j-1}/P_{j-1,m}+13rc
                     \qquad(0\le m<r).
 \end{split}
\end{equation}
\end{paramdefinition}
The family of test polynomials is indexed by primitive vectors of the true lattice. The two signs of the diagonal unit are treated by Lemma~\ref{lem:plane-digit-units-and-cuts}. The reserves are inequalities on the integer conductor chosen for the small-prime calculation; its specialization is made in Section~\ref{sec:host-edge-tests}.


\begin{proposition}[An integral polynomial with prescribed orders at the model cuts]\label{prop:plane-corrected-gram-polynomial}
Use the lattices, orders and sign approximations in
Definition~\ref{hyp:plane-gram-data}, including \(\nu(F)=0\)
and the reserves \eqref{eq:plane-gram-reserves}.
Then there is an integral model polynomial \(Z_*\) of collapsed
degree \(d\), with model link orders at least
\(\ell_{jm}-30c\) and chain orders at least
\(\theta_{ji}-s_{ji}-12c\), such that
\[
 v([F,\gamma\phi Z_*])-n_0+K_d
                               \le D_l+D_m+4K_d.
\]
The order and valuation bounds are uniform over the coefficients,
digit classes and radices with the data of Definition~\ref{hyp:plane-gram-data}.
\end{proposition}

\begin{proof}
\proofstep{Projecting a monomial to the link cuts.}
Take a monomial \(x^\alpha\) attaining \(\nu(F)=0\) and
project it to the true link cuts at levels \(\ell_{jm}\).
The result \(Y\) has \(\nu(Y)\ge h_\alpha-D_l\), and
\[
 [F,Y]=F_\alpha\,\alpha!
\]
has valuation \(h_\alpha\).
Apply an integral binomial approximation \(g\) of \(\gamma^{-1}\)
of degree at most \(6c\) to its coefficients. A factor
\(\binom{\alpha_x}{a}\) acts as
\(x^a\partial_x^{[a]}\). Its loss in a signed cut is at
most \(a\), by the cut estimates for the derivative and
multiplication. This bounds the total link loss by \(6c\).
The change in the pairing from \([F,Y]\) to
\([F,\gamma(gY)]\) has valuation at least
\[
 c+h_\alpha-D_l>h_\alpha.
\]

Ordinary coefficients of \(gY\) have valuation at least
\(h_\alpha-D_l-n_0\). Multiplication by
\[
 p^{\lceil n_0+D_l-h_\alpha\rceil}
\]
therefore makes \(gY\)
an integral polynomial \(G\) in the original variables. It has link orders at least
\(\ell_{jm}-6c\) and its twisted pairing has valuation at
most \(n_0+D_l+1\).

\proofstep{Lifting \(G\) in the prescribed plane.}
To lift \(G\), first put each \(X_t\) entirely in \(L_t\).
Within frame \(j\), order its block letters first and then
the shared letters in decreasing order of \(m\). Each
\(A_{jm}\) is a prefix of this order. Use the integral basis
of signed path words proved in Section~\ref{sec:column-lifts},
now testing the least torus weight at the prefixes.
In a positive column the least exponent selects a
subset avoiding the torus coordinate whenever \(h<r\).
In a nonpositive column it selects a dual subset
containing the coordinate whenever \(h<r\). Thus if
\(h_b\) are the prefix heights,
\[
 \lambda_r=
 \sum_{b\ {\rm positive}}\mathbf1_{h_b=r}
       -\sum_{b\ {\rm nonpositive}}\mathbf1_{h_b<r}.
\]
The Pieri decomposition at the successive prefixes
selects exactly the paths whose displayed sums obey all
the lower bounds. The associated words span the
intersection of the cuts. Since they are a
subset of the integral basis, the expansion of an integral
polynomial in this intersection uses integral coefficients.

In a path word, a column meeting at least one block
letter moves all its shared occurrences with \(m>0\) from
\(L_{j,m}\) to \(T_{j,m}\). Keep \(L_{j,0}\) and every
occurrence in a column disjoint from the blocks. Each
column uses a given shared letter at most once. Count the columns by one of their block occurrences. Each chosen occurrence lies in only one column, so
\[
 \#\{\text{moved columns}\}\le\#\{\text{block occurrences}\}\le b_j.
\]
This gives the tangent degree bounds.
On a moved column the new cut \(B_j\cup\{T_{j,m'}:m'>m\}\)
acts on exactly the old cut occurrences. An unmoved
column has zero order there. Its contribution to the
prescribed lower bound was at most zero.

In fact for a
positive column the height at the cut is bounded by the
number of its occupied vector letters before the cut,
which is \(<r\) since the column meets only shared letters.
For a prefix, a moved column retains its signed cut
order and an unmoved column contributes zero, at least
its previous order. Minimum Laurent orders add on products, giving the required cut
bound for the whole word. For the coefficient valuation the relevant
identity is
\[
      \nu_{\rm coeff}(AB)=\nu_{\rm coeff}(A)+\nu_{\rm coeff}(B)
                 \qquad(A,B\ne0).
\]
Indeed, divide each factor by its minimum uniformizer power. Their
reductions are nonzero polynomials over the residue field, so their
product has nonzero reduction. The injective renaming on each
column preserves that column's \(\nu_{\rm coeff}\), and
hence preserves the valuation of the product word.
Collapsing \(T_{j,m}\) back to the corresponding shared
letter restores each column separately.

Perform the operation coefficientwise outside each frame.
Since the renamed letter sets are disjoint, that operation
commutes with the cut substitutions in the other frames.
For a remaining shared variable, divide its exponent by its radix.
The lift and collapse act on the monomial by
\[
 a_t=P_tq_t+r_t,\quad 0\le r_t<P_t,\qquad
 L_t^{r_t}W_t^{q_t}\:\longmapsto\: X_t^{r_t+P_tq_t}=X_t^{a_t}.
\]
The tangent exponents stay fixed throughout. To see that it preserves
the link bounds, group by the \(L\)-exponents. These letters
are outside the new link cuts, so coefficient extraction
preserves the bounds. The added monomial in the \(W\) letters
has nonnegative order at every model link cut, since those
letters are vectors. The integral model polynomial \(Z\)
therefore has admissible support, \(\phi(Z)=G\), and
\(\operatorname{ord}_{U_{jm}}Z\ge\ell_{jm}-6c\).

\proofstep{Correction to impose the chain cuts.}
For an admissible model exponent \(\zeta\), let \(A\) be
its collapsed exponent without \(T\), and let \(M\) be its
tangent exponent on the true shared letters. Thus
\(\phi(x^\zeta)=x^{A+M}\). The digit factorial identity
gives
\[
       A! M! = \kappa_B\,u_\zeta\,\zeta!,
       \qquad
       \eta_\zeta=\gamma(A)u_\zeta ,
\]
where \(u_\zeta\) is the product of the digit units.
The constant \(\kappa_B\) depends on the individual label and
tangent exponents and on the model multidegrees, which we fix
in each block of the next projection. Explicitly it is the
product of the digit constants for each \(W_t\)-degree divided
by the product of the label factorials. The letter \(T_t\)
contributes the same factorial on both sides.

\proofstep{Diagonal correction factors.}
Choose diagonal multipliers \(g_+\) and \(g_-\) approximating
\(\eta_\zeta\) and its inverse modulo \(p^c\), with
integral values and cut cost \(12c\) on cuts excluding \(L\).
For fixed label exponents, substitute
\(\rho_{t,a}+P_t\beta_{t,a}\) into the binomial approximations
to \(\gamma^{\pm1}\). Apply forward differences in \(\beta\) to the integer-valued polynomial \(A(\beta)=\binom{\rho+P\beta}{k}\). Its expansion is
\[
 A(\beta)=\sum_{j=0}^k(\Delta^j A)(0)\binom{\beta}{j},\qquad (\Delta^j A)(0)\in\mathbb Z.
\]

Multiplying this approximation by the degree-\(6c\) digit
unit approximation yields the cost \(12c\). Extracting
coefficients of \(L\) leaves every such cut intact, so the
estimate holds after adding the label classes.

In each model multidegree let \(\Pi_m\) be the product
projector on chain orders at least \(\theta_{ji}-s_{ji}\).
Use its linear extension over all multidegrees. It fixes
the exponents on \(L,T\), preserves the admissible support
conditions, and commutes with all model link projections.
Each multidegree component of a polynomial satisfying a
link bound still satisfies it, by the homogeneous cut action.
 Set
\[
 \Omega=g_-\Pi_m g_+.
\]

Its chain orders are \(\ge\theta_{ji}-s_{ji}-12c\).
Any model link bound is lowered by at most \(24c\).
For \(\zeta,\zeta'\) in the same fixed block of \(\Pi_m\),
\[
 \lvert v(\zeta!)-v(\zeta'!)\rvert\le K_d.
\]

Only the block and \(W\) exponents can vary. Their degrees
are bounded by the corresponding true degrees, so the
multinomial estimate applies without a second factor for
the fixed \(L,T\) exponents. The factorial-weighted projection
bound proves that \(p^{N_2}\Omega\) preserves integrality
on admissible support, with
\(N_2=\lceil D_m+K_d/2\rceil\).

\proofstep{Precision of the corrected pairing.}
For a fixed tangent exponent \(M\) put
\[
         F_M=\gamma^{-1}\,\partial^{[M]}(\gamma F).
\]
Every coefficient of \(F_M\) has valuation \(\ge -n_0\):
the unit multipliers and the integral divided derivative
preserve the original coefficient bound. At a fixed label
\(\rho\), its digit part has a decomposition \(Q+E\) such
that
\[
 v(E_\beta)\ge c-n_0
\]
and
\(\operatorname{ord}_{S_{ji}}Q\ge\theta_{ji}-s_{ji}\)
for all \(j,i\). To obtain it, start with the digit parts
of \(F\). In each of the two twist operations replace the
unit multiplier by its degree-\(6c\) approximation. Each
replacement error has coefficient valuation at least
\(c-n_0\), and each approximation loses at most \(6c\) in
the chain orders, using the affine binomial expansion.
Between the two multipliers use the divided derivative
expansion of the preceding subsection on each shared letter.

The divided-derivative expansion takes the input label to the residue of \(\rho-M\)
coordinatewise. The error coefficients have valuation at
least \(c-n_0\). The approximate derivatives and all
multiplying monomials have integral coefficients, so errors
already made retain this valuation. Only derivatives on
the \(W\) letters in the cut can lower its order. Their
total loss is bounded by
\[
    \sum_{m=0}^{r-1}
       \left(\frac{|M_{j-1,m}|}{P_{j-1,m}}+c\right)
    \ \le\ s_{ji}-12c
\]
by \eqref{eq:plane-gram-reserves} and \(r\ge1\).
This produces \(Q,E\) with the asserted relation.

\proofstep{Comparing factorial terms in the corrected pairing.}
For an admissible monomial the factorial identity and the
coefficient formula for the divided derivative give
\[
 [F,\gamma\,\phi(x^\zeta)]
       = \kappa_B\,\eta_\zeta\, (F_M)_A\,\zeta! .
\]
In the model multidegree of \(\zeta\), multiply \(Q\) by
the fixed label and tangent monomial, then keep that
multidegree. The resulting polynomial \(Q_B\) still has
the required chain orders because \(L,T\) are outside
those cuts. On that block the difference
\[
    g_-(\zeta')\,\eta_{\zeta'}(F_M)_{A'}-(Q_B)_{\zeta'}
\]
has valuation at least \(c-n_0\). The pairing with
\(Q_B\) is unchanged by \(\Pi_m\).

To estimate the remaining
error from a coefficient of \(\Pi_m x^\zeta\), use its bound
\(h_\zeta-h_{\zeta'}-D_m\). From the split factorial
bound applied to \(\zeta\) and \(\zeta'\), and the equality
of their \(\kappa_B\), we have
\[
 v(\kappa_B)+h_\zeta+h_{\zeta'}\ \ge\ 2n_0-2K_d .
\]
The error terms in comparing the evaluation before and
after \(\Pi_m\) thus have valuation at least
\(c+n_0-2K_d-D_m\), including the original monomial term.
Applying \(g_+\) at the input preserves this error estimate.
The factor \(g_+(\zeta)g_-(\zeta)-1\) has valuation at least
\(c\), and the original evaluation has valuation at least
\(n_0-K_d/2\) by \(\nu(F)=0\) and the true factorial bound.
Linearity and integral coefficients of \(Z\) now give
\begin{equation}\label{eq:plane-gram-stability}
 v\bigl([F,\gamma\phi(\Omega Z)]-[F,\gamma\phi(Z)]\bigr)
               \ \ge\ c+n_0-2K_d-D_m .
\end{equation}

\proofstep{The available precision after the conductor loss.}
The strict conductor inequality in \eqref{eq:plane-gram-reserves}
makes \eqref{eq:plane-gram-stability} larger than
\(n_0+D_l+1\). Hence the integral test
\(Z_*=p^{N_2}\Omega Z\) has valuation of its pairing at
most \(n_0+D_l+1+N_2\). Its model link and chain orders are
\(\ge\ell_{jm}-30c\) and \(\ge\theta_{ji}-s_{ji}-12c\),
respectively. Its collapse has true multidegree \(d\) and
integral coefficients, hence \(\nu(\phi Z_*)\ge n_0-K_d\).
In particular
\[
 v([F,\gamma\phi Z_*])-n_0+K_d
                 \le D_l+D_m+4K_d .
\]
Here \(K_d\ge1\) absorbs the two roundings. The estimate is uniform and needs no error polynomial.
The chain and link losses and the valuation gain are
uniform over all input coefficients, label classes, and
radices at the specified degrees and cuts.
\end{proof}
